\documentclass[11pt, a4paper]{article}
\usepackage[utf8]{inputenc}
\usepackage{kotex}
\usepackage[margin=20mm]{geometry}
\usepackage{graphicx}
\usepackage{booktabs}
\usepackage{xcolor}
\usepackage{float}
\usepackage{enumitem}
\usepackage{fancyhdr}
\usepackage{array}
\usepackage{colortbl}
\usepackage{tcolorbox}

\definecolor{gapred}{HTML}{E53935}
\definecolor{gapblue}{HTML}{1565C0}
\definecolor{headerblue}{HTML}{1A237E}
\definecolor{lightgray}{HTML}{F5F5F5}
\definecolor{boxblue}{HTML}{E3F2FD}

\pagestyle{fancy}
\fancyhf{}
\lhead{\small\color{gray} AI 정책 공론-입법 격차 분석}
\rhead{\small\color{gray} 2026.04}
\cfoot{\thepage}

\setlength{\parindent}{0pt}
\setlength{\parskip}{6pt}

\begin{document}

% ── 표지 ──
\begin{center}
\vspace*{30mm}
{\color{headerblue}\rule{\textwidth}{2pt}}
\vspace{8mm}

{\Huge\bfseries AI 공론과 입법의 격차}\\[4mm]
{\Large 한국·미국·EU 비교 분석 보고서}

\vspace{10mm}
{\large 2026년 4월}

\vspace{8mm}
{\color{headerblue}\rule{\textwidth}{2pt}}

\vspace{40mm}
\begin{tcolorbox}[colback=boxblue, colframe=headerblue!50, width=0.85\textwidth]
\textbf{핵심 발견}\\[2mm]
\textbf{시장경쟁/독과점(Market/Antitrust)} 이슈가 한국·미국·EU \textbf{3국 모두}에서 가장 큰 입법 사각지대로 확인됨. 뉴스에서 15--29\% 비중이나 관련 법안은 0--4\%에 불과.
\end{tcolorbox}
\end{center}

\newpage

% ── 1. 분석 개요 ──
\section*{\color{headerblue}1. 분석 개요}

\begin{tcolorbox}[colback=lightgray, colframe=gray!50]
\textbf{목적} \quad AI에 대한 공적 담론(뉴스)과 실제 입법 활동 사이의 격차를 정량 측정\\
\textbf{방법} \quad GPT-4.1-mini 기반 10개 정책 속성 자동 분류\\
\textbf{기간} \quad 2016.03 -- 2026.04 (뉴스) / 각국 현 회기 (법안)
\end{tcolorbox}

\vspace{4mm}
\textbf{분석 데이터 규모}

\begin{table}[H]
\centering
\begin{tabular}{lrl}
\toprule
\textbf{구분} & \textbf{건수} & \textbf{출처} \\
\midrule
뉴스 (영국) & 8,104건 & Guardian \\
뉴스 (미국) & 2,797건 & New York Times \\
뉴스 (한국) & 14,328건 & 네이버 뉴스 \\
\midrule
법안 (미국) & 207건 & 118th + 119th Congress \\
법안 (한국) & 121건 & 22대 국회 \\
법안 (EU) & 884건 & AI Act 조문 + 수정안 \\
\bottomrule
\end{tabular}
\end{table}

% ── 2. 주요 발견 ──
\section*{\color{headerblue}2. 주요 발견}

\subsection*{발견 1: 시장경쟁/독과점 — 3국 공통 입법 사각지대}

\begin{figure}[H]
\centering
\includegraphics[width=\textwidth]{figures/fig5_cross_country.pdf}
\caption*{[그림 1] 국가별 공론-입법 격차 비교}
\end{figure}

시장경쟁/독과점(Market/Antitrust) 관련 뉴스 보도 비중은 15--29\%에 달하나, 관련 법안은 한국 0\%, 미국 3--4\%, EU 1--2\%에 불과함. \textbf{3국 모두에서 가장 큰 격차}를 보이는 공통 입법 사각지대임.

\subsection*{발견 2: 한국 — 산업정책 편중, 공익법안 과잉}

\begin{figure}[H]
\centering
\includegraphics[width=\textwidth]{figures/fig2_gap_korea.pdf}
\caption*{[그림 2] 한국 공론-입법 격차}
\end{figure}

\begin{itemize}[noitemsep]
    \item 네이버 뉴스의 \textbf{47.6\%}가 산업정책 관련 — ``육성 프레임'' 지배적
    \item 시장경쟁/독과점 관련 법안 \textbf{0건} — 가장 심각한 공백
    \item 공익/소비자보호 법안이 27.3\%로 뉴스(5.8\%) 대비 \textbf{과잉 집중}
\end{itemize}

\subsection*{발견 3: 미국 — 정권 교체에 따른 입법 기조 급변}

\begin{figure}[H]
\centering
\includegraphics[width=\textwidth]{figures/fig3_gap_us.pdf}
\caption*{[그림 3] 미국 118th vs 119th Congress 비교}
\end{figure}

\begin{itemize}[noitemsep]
    \item 118th(바이든): 책임/윤리AI 법안 \textbf{27.9\%} $\to$ 119th(트럼프): \textbf{3.8\%}로 급감
    \item 119th: Safety \textbf{37.7\%}, 국가안보 \textbf{28.3\%}로 급증
    \item Market/Antitrust 공백은 양 회기 모두 지속 (+16\%p 이상)
\end{itemize}

\subsection*{발견 4: EU — Safety 단일 목적 입법}

EU AI Act 113개 조항 중 \textbf{74\%가 Safety} 관련. 수정안 771건에서도 65\%가 Safety에 집중. 노동(0\%), 저작권(0\%), 선거(0\%)는 별도 법률로 처리.

% ── 3. 비교표 ──
\section*{\color{headerblue}3. 정책 속성별 비교표}

\begin{figure}[H]
\centering
\includegraphics[width=\textwidth]{figures/fig1_heatmap.pdf}
\caption*{[그림 4] 8개 데이터셋의 정책 속성별 비중(\%) 히트맵}
\end{figure}

% ── 4. 시사점 ──
\section*{\color{headerblue}4. 시사점}

\begin{tcolorbox}[colback=boxblue, colframe=headerblue!50]
\textbf{입법자 대상}
\begin{itemize}[noitemsep]
    \item 시장경쟁/독과점, 노동, 저작권 영역의 입법 공백 우선 해소 필요
    \item 산업정책 외 다양한 AI 정책 속성에 대한 균형 있는 입법 필요
\end{itemize}

\textbf{정책입안자 대상}
\begin{itemize}[noitemsep]
    \item 국가간 비교를 통한 AI 정책 포트폴리오 벤치마킹
    \item 정권 교체에 영향받지 않는 초당파적 AI 거버넌스 체계 필요
\end{itemize}
\end{tcolorbox}

% ── 5. 한계 ──
\section*{\color{headerblue}5. 한계 및 향후 과제}
\begin{itemize}[noitemsep]
    \item 네이버 뉴스 수집 기간 제한 (6주) — BigKinds 활용 확장 예정
    \item Guardian/NYT 중도좌파 편향 — 보수 매체 무료 API 부재
    \item GPT 분류의 해석 가능성 제한 — temperature=0으로 재현성 확보
    \item 법안 수 기반 분석의 한계 — 법안 내용의 질적 차이 미반영
\end{itemize}

\vfill
\begin{center}
\small\color{gray}
본 보고서의 데이터 수집·분류·시각화 전 과정은 스크립트로 구현되어 재현 가능합니다.
\end{center}

\end{document}
